\documentclass[10pt]{article}
\usepackage[margin=20mm]{geometry}
\usepackage{amsmath,amssymb,fontspec,unicode-math,xurl,hyperref}
\IfFontExistsTF{DejaVu Serif}{\setmainfont{DejaVu Serif}}{\setmainfont{Latin Modern Roman}}
\setmathfont{latinmodern-math.otf}
\setlength{\parindent}{0pt}\setlength{\parskip}{6pt}\sloppy
\title{Technical Verification of Transformer Training Dynamics for Regular Language Recognition}
\author{AI-assisted technical review}
\date{}
\begin{document}
\maketitle
Abstract: The Phase 2 proof neglects the scaling dependence of its max-margin norm. Under the stated Phase 1 updates that norm grows at least linearly with the attention scale, invalidating the chosen proof horizon and placing a later norm threshold beyond the interval where the needed lemma is established. This core derivation failure supports Major. A separate two-step counterexample to query-position symmetry supports Minor. Neither finding asserts empirical failure.\par
\section*{Identity and scope}
This report reviews Ruiquan Huang, Yingbin Liang, and Jing Yang, How Transformers Learn Regular Language Recognition: A Theoretical Study on Training Dynamics and Implicit Bias, ICML 2025, PMLR 267, pp. 25516–25549 (I25-43533), in the exact staged published-proceedings PDF. Page numbers mean printed internal pages 1–34, equal to the one-based PDF index.\par
The main training formulation, Phase 1 and Phase 2 claims, mathematical appendices, and experiment descriptions were examined. The independent check independently checked the operative training rule, assumptions, margin definition, horizon and proof dependencies, and the two-step counterexample. Source pages 3, 4, 6, 13, 23, 26, and 28 were visually inspected. One central derivation problem supports Major; a separate local false symmetry statement supports Minor. These findings do not establish that the tasks cannot be learned or that the experiments are false.\par
\section*{Model and assumptions}
For every binary string \(s\) of length \(L\le L_{\max}\), write \(X_s\) for its matrix of orthonormal position–token embeddings, and \(x_{s,L}\) for its last token. The attention representation and scalar prediction are \[v_s(W)=X_s\operatorname{softmax}(X_s^\top W x_{s,L}/\lambda),
\qquad f_s(u,W)=\langle u,v_s(W)\rangle.\] The even-pairs label is \(y_s=1\) if the first and last tokens agree, and \(-1\) otherwise. The loss sums the average logistic loss at each length: \[\mathcal L(u,W)=\sum_{L=1}^{L_{\max}}2^{-L}
\sum_{s\in\{a,b\}^L}\log(1+e^{-y_sf_s(u,W)}).\] Initialization is \(u_0=W_0=0\). In Phase 1, \[u_{t+1}=u_t-\eta\nabla_u\mathcal L_t,\qquad
W_{t+1}=W_t-\eta\lambda\nabla_W\mathcal L_t.\] Theorem 4.1 (p. 5) chooses \(\lambda=\Omega(L_{\max}^2)\), \(t_0=O((\eta L_{\max})^{-1})\), and \(\eta=O(\min\{L_{\max}^{-1},\lambda^{-2/3}\})\). Phase 2 uses learning rate \(\eta\) for both parameter blocks. The max-margin vector \(u^\ast\) is the minimum-norm vector satisfying \(y_s\langle u^\ast,v_s(W_{t_0})\rangle\ge1\) for every training string. The following analysis grants the claimed Phase 1 separability.\par
\section*{Substantive findings}
Problem 1 (Incompatible margin dependence and Phase 2 proof horizon).  The horizon used to prove the central Phase 2 claims is inconsistent with the margin dependence generated by the stated Phase 1 dynamics.\par
\subsection*{Original claims and proof dependency}
Equation reference: eq:T, eq:t1. Theorems 4.4–4.5 (p. 6) give alignment and loss bounds up to a horizon of order \(\lambda^{2/3}/(\eta L_{\max})\). Appendix B’s parameter setup (p. 23) explicitly writes \[
T=\min\left\{
\frac{\lambda}{\sqrt{3\|u^\ast\|}\,\eta L_{\max}},
\frac{\lambda^{2/3}}{2^{2/3}\eta L_{\max}}
\right\}
=\frac{\lambda^{2/3}}{2^{2/3}\eta L_{\max}}.\] Lemma B.5 is valid only for \(t_0<t\le T\). The proof of Lemma B.6 (p. 26) then introduces \[
t_1=t_0+1+\frac{14\|u^\ast\|^2}{\eta}
\exp\left(\frac{5}{3\Delta\|u^\ast\|}\right),\qquad \Delta>0,\] and invokes Lemma B.5 for \(t\ge t_1\). This is used in Theorem B.7 (p. 28), restating the loss result. Appendix C (pp. 28, 33) explicitly reuses the Phase 2 argument for parity training.\par
\subsection*{Complete check of the parameter dependence}
Equation reference: eq:Wbound. The logistic derivative has absolute value at most one, \(\|v_s(W)\|\le1\), and the softmax Jacobian has operator norm at most one. Since the loss has one average term per length, \[\|\nabla_u\mathcal L_t\|\le L_{\max},\qquad
\|\nabla_W\mathcal L_t\|\le L_{\max}\|u_t\|/\lambda.\] Thus \(\|u_t\|\le\eta L_{\max}t\), and the actual Phase 1 update gives \[
\|W_{t_0}\|\le
\eta L_{\max}\sum_{s=0}^{t_0-1}\|u_s\|
\le \frac{\eta^2L_{\max}^2t_0(t_0-1)}2=:B.\] Write \(t_0\le c/(\eta L_{\max})\) with the fixed constant allowed by the stated asymptotic regime. Then \(B\le c^2/2\), independently of \(\lambda\).\par
Equation reference: eq:margin. Now retain only the four length-two strings \(aa,bb,ab,ba\), with labels \(+,+,-,-\). At uniform attention, \[v_{aa}(0)+v_{bb}(0)-v_{ab}(0)-v_{ba}(0)=0.\] Orthonormal columns and the softmax Lipschitz bound imply \(\|v_s(W)-v_s(0)\|\le\|W\|/\lambda\). For any feasible unit-margin separator, summing its four margin inequalities yields \[4\le\left\langle u,\sum_s y_s v_s(W)\right\rangle
=\left\langle u,\sum_s y_s(v_s(W)-v_s(0))\right\rangle
\le4\|u\|\|W\|/\lambda.\] Consequently, whenever the promised separator exists, \[
\|u^\ast\|\ge\lambda/\|W_{t_0}\|
\ge\lambda/B\ge2\lambda/c^2.\] If \(W_{t_0}=0\), these four features are not separable; this cannot rescue the asserted Phase 2 argument.\par
For the second entry to be the minimum in [eq:T] (\#eq:T), one needs \[\|u^\ast\|\le \frac{2^{4/3}}{3}\lambda^{2/3}.\] This contradicts [eq:margin] (\#eq:margin) as \(\lambda\) grows with fixed \(c\) and \(L_{\max}\ge2\). The condition \(\lambda=\Omega(L_{\max}^2)\) does not provide the missing upper bound on the margin norm. An admissible scaling family is \(L_{\max}=2\), \(\lambda=n\), \(\eta=1/n\), and \(t_0=\lfloor cn/2\rfloor\), with fixed sufficiently small \(c>0\); it satisfies the stated step-size and Phase 1 orders.\par
Furthermore, [eq:t1] (\#eq:t1) implies \[t_1>\frac{14\|u^\ast\|^2}{\eta}
\ge\frac{56\lambda^2}{c^4\eta},
\qquad
\frac{t_1}{T}>
\frac{56\,2^{2/3}L_{\max}}{c^4}\lambda^{4/3}.\] The proof’s chosen sufficient threshold therefore exceeds the proved horizon in that regime. Restoring the literal minimum in [eq:T] (\#eq:T) only makes the horizon smaller. The proof has not established a nonempty time interval where all its premises apply.\par
\subsection*{Finite diagnostic and loss implication}
For \(\eta=10^{-3}\), \(\lambda=100\), \(L_{\max}=2\), and \(t_0=50\), the independent calculation gives \(B=0.0049\), \(\|u^\ast\|\ge20408.1633\), second horizon \(6786.0440\), first horizon at most \(202.0726\), and \(t_1>5.8309\times10^{12}\). These are evaluations of analytic bounds, not a claimed numerical solution of the max-margin problem. The asymptotic argument is the decisive contradiction and does not depend on choosing unspecified constants for this illustration.\par
The displayed loss bound is \(\mathcal L_t=O(L_{\max}\|u^\ast\|^2/(\eta\sqrt t))\). Substitution of \(T=C\lambda^{2/3}/(\eta L_{\max})\) gives the expression \[O\!\left(
\frac{L_{\max}^{3/2}\|u^\ast\|^2}
{\sqrt{C\eta}\lambda^{1/3}}
\right),\] not the displayed positive half-power of \(\eta\) on pp. 6 and 28. The margin norm also cannot be held fixed while varying \(\lambda\). An increasing upper-bound expression is not a lower bound on the actual loss; it shows that this argument does not justify arbitrary accuracy.\par
\subsection*{Impact and counterevidence}
This is Major: a demonstrated invalid derivation directly supports the central alignment, loss-convergence, and reused parity guarantees. Phase 1 separability was granted, and the experiments and qualitative two-phase behavior are not refuted. Hidden constants that depend arbitrarily on \(\lambda\) cannot justify a conclusion obtained by sending \(\lambda\) upward. A replacement result needs a proved compatible time regime and complete margin dependence. No such complete repair is certified.\par
Problem 2 (Query-position symmetry in Lemma B.1 is false).  Lemma B.1 (p. 13) asserts \[\langle E_\ell^w,W_tE_L^w\rangle
=\langle E_\ell^{-w},W_tE_L^w\rangle
\quad\text{for all }\ell\ge2.\] Its sample-pairing argument incorrectly treats flipping the last token as label-preserving.\par
\subsection*{Admissible two-step counterexample}
Set \(L_{\max}=2\), use the full six-string training set with the stated length weights and orthonormal embeddings, and start from zero. Choose positive \(\eta,\lambda\) in the stated regime with two Phase 1 updates. Direct differentiation gives \[u_1=\frac{\eta}{4}(E_1^a+E_1^b),\qquad W_1=0.\] At the next update, \[\langle E_2^a,W_2E_2^a\rangle
=\frac{\eta^2}{64}\tanh(\eta/16)>0,\qquad
\langle E_2^b,W_2E_2^a\rangle=0.\] Indeed, the latter coordinate never appears with query \(E_2^a\): a length-two sequence with last token \(a\) has no position-two token \(b\). The two same-query sequences produce the stated nonzero diagonal after their logistic contributions are combined.\par
The independently written enumeration at \(\eta=0.001,\lambda=100\) returns \(9.76562498728945\times10^{-13}\) versus zero; the analytic diagonal is \(9.765624987284342\times10^{-13}\). The tiny rounding difference between the two nonzero evaluations does not affect the strict inequality.\par
\subsection*{Bounded impact}
This is Minor. The false equality concerns a query-position coordinate that is inactive for that query. Interior-position and global token-flip symmetries may still support the intended sample reduction; no additional core failure is established from this defect. Correcting the position domain is necessary, but a boundary edit has not been certified to repair every downstream proof. This judgment is based on functional dependence, not appendix location.\par
\section*{Conclusion and remaining work}
Issue | Assessment | Repair verification | Author adoption | Remaining work\par
P1 | Major | No replacement theorem | Unknown | Prove a compatible horizon with margin dependence\par
P2 | Minor | Boundary restriction proposed | Unknown | Check the corrected induction and downstream uses\par
Scientific adjudication is pending for every issue; human reviews remain empty. Author adoption has not been requested.\par
The findings apply to the frozen PMLR PDF. Cached arXiv metadata was inspected by the auxiliary, but no alternate full-text revision was compared. No absence of an erratum is claimed. Empirical training was not reproduced, external cited theorems were not independently re-proved, and not every Phase 1 induction constant was checked.\par
9 R. Huang, Y. Liang, and J. Yang, How Transformers Learn Regular Language Recognition: A Theoretical Study on Training Dynamics and Implicit Bias, ICML, PMLR 267 (2025), 25516–25549. Published PDF with appendices. https://proceedings.mlr.press/v267/huang25s.html.\par
\section*{Supporting checks}
The Phase 2 argument was independently challenged, with original pp. 6, 23 and 26 checked. The gap remains, while \(t_1\) is the proof’s chosen sufficient threshold, not a necessary hitting time.\par
\end{document}
